\subsection{Purpose}
The increasing focus on sustainability and eco-friendly mobility has encouraged many people to use bikes to move around the city, both for daily trips and for sport.
However, traditional navigation systems are not designed specifically for cyclists. They are mainly optimized for the shortest distance or travel time and do not show important information for a bike trip, such as road surface quality, typical car speed on those streets, or the presence of obstacles and potholes, which are essential for a cyclist who wants to plan a safe and comfortable route. Cyclists, especially those with sports goals, also need a way to keep an inventory of their past routes and progress.
The purpose of Best Bike Paths (BBP) is to provide cyclists a place where they can find, rate, and share good bike paths and manage their own. 
The system records trips, collects sensor and weather condition data, and lets users add details about the status and of paths. Users can publish their tracks, explore and rate routes made by others, and in this way a community of cyclists is created. Among the shared routes, the application highlights and suggests those that are considered the safest and best rated by the community.In this way, BBP helps make cycling safer, more informed, and more enjoyable for its community.
\subsubsection{Goals}
\textbf{G1:} Users can manually create bike paths or alternatively be assisted by the system to automatically track the paths thanks to the mobile device's sensors. Subsequently, the paths can be reviewed and enriched with general information and with a status before being saved inside the inventory. \\ [2mm]
\textbf{G2:} Users can visualize their own inventory and search for their saved paths with their statistics and status that are enriched by meteorological information provided by external services. \\ [2mm] 
\textbf{G3:} Users can publish their paths, sharing them with the community, creating a social component in the system where users can discover different paths proposed by others. \\ [2mm]
\textbf{G4:} Given a starting point and a destination, the software suggests the best route to the user considering a ranking created from the scores assigned to each path and evaluated based on the status provided by one or more users and its effectiveness.
\subsection{Scope}
The core function of BBP is to support cyclists in recording, managing, and sharing bike paths, and to help all users find safe and suitable routes between two locations. The final user of the software is a cyclist that interacts with BBP through a personal account (not mandatory) on a mobile application. They can create bike paths manually or let the system track them automatically using GPS and  sensors on their mobile device. After a trip, they can review the recorded path, confirm or correct detected problems, and enrich it with information on status and other details before saving it in their personal inventory. BBP stores each path together with statistics such as distance, time, average speed, and enriches these trips with meteorological data obtained from external weather services. Registered users can publish some of their paths, making them visible to the community. Cyclists view public paths, see their status and obstacles, and discover paths shared by the community. In this way, the system introduces a social component where cyclists can explore and rate different paths. Given a starting point and destination, BBP searches its inventory of recorded and published paths and suggests routes to the user. The system ranks the candidate routes using scores based on their status and the feedback of the users. BBP focuses on collecting, aggregating, and presenting information about bike paths among a community of users.
\subsubsection{World Phenomena}
\subsubsection{Shared Phenomena}
\subsection{Definitions, Acronyms, Abbreviations}
\subsubsection{Definitions}
\subsubsection{Acronyms}
\subsubsection{Abbreviations}
\subsection{Revision History}
\subsection{Reference Documents}
\subsection{Document Structure}